\section{Неоднородные системы. Теорема об общем решении неоднородной системы}

\begin{definition}
Система линейных уравнений $Ax = b$ называется \textbf{неоднородной}, если вектор свободных членов $b \neq 0$. Соответствующая ей система $Ax = 0$ называется \textbf{однородной}.
\end{definition}

\begin{theorem}[Об общем решении неоднородной системы]
Пусть неоднородная система $Ax = b$ совместна. Пусть $X_{ч}$ --- какое-либо её частное решение, а $X_{оо}$ --- общее решение соответствующей однородной системы $Ax = 0$. Тогда общее решение неоднородной системы имеет вид:
\[
X = X_{оо} + X_{ч}.
\]
\end{theorem}

\begin{proof}
Доказательство состоит из двух частей.

\begin{enumerate}
    \item \textbf{Достаточность:} Покажем, что вектор вида $X = X_{оо} + X_{ч}$ является решением неоднородной системы $Ax = b$.
    Подставим этот вектор в систему и воспользуемся свойством линейности матричного умножения:
    \[
    A(X_{оо} + X_{ч}) = A X_{оо} + A X_{ч}.
    \]
    По определению, $X_{оо}$ --- решение однородной системы, поэтому $A X_{оо} = 0$. А $X_{ч}$ --- частное решение неоднородной системы, поэтому $A X_{ч} = b$. Следовательно:
    \[
    A(X_{оо} + X_{ч}) = 0 + b = b.
    \]
    Таким образом, $X_{оо} + X_{ч}$ действительно является решением системы $Ax = b$.

    \item \textbf{Необходимость:} Покажем, что любое решение $X$ неоднородной системы может быть представлено в таком виде.
    Пусть $X$ --- произвольное решение системы $Ax = b$, то есть $AX = b$. Рассмотрим вектор $y = X - X_{ч}$. Проверим, чему равно $Ay$:
    \[
    Ay = A(X - X_{ч}) = AX - AX_{ч}.
    \]
    Так как и $X$, и $X_{ч}$ являются решениями неоднородной системы, то $AX = b$ и $AX_{ч} = b$. Подставляем:
    \[
    Ay = b - b = 0.
    \]
    Это означает, что вектор $y$ является решением однородной системы $Ax = 0$. Следовательно, $y$ может быть представлен как часть общего решения однородной системы, то есть $y = X_{оо}$. 
    Возвращаясь к обозначению $y = X - X_{ч}$, получаем:
    \[
    X_{оо} = X - X_{ч} \quad \Rightarrow \quad X = X_{оо} + X_{ч}.
    \]
\end{enumerate}
Теорема полностью доказана.
\hfill$\blacksquare$
\end{proof}

\begin{corollary}
Если $\operatorname{rang} A = \operatorname{rang} \overline{A} = r < n$, то общее решение неоднородной системы содержит ровно $n - r$ произвольных постоянных и записывается как сумма конкретного частного решения и линейной комбинации векторов ФСР соответствующей однородной системы.
\end{corollary}
